% ============================================================
%  C: Como Programar -- 6ª Edição | Deitel & Deitel
%  Capítulo 2 -- Introdução à Programação em C
%  EXERCÍCIOS DE AUTORREVISÃO (2.1 e 2.2) -- com respostas
%  Abordagem incremental: Caps. 1 + 2
% ============================================================
\documentclass[12pt,a4paper]{article}
\usepackage[utf8]{inputenc}
\usepackage[T1]{fontenc}
\usepackage[brazil]{babel}
\usepackage{geometry}
\geometry{top=2.5cm,bottom=2.5cm,left=3cm,right=3cm}
\usepackage{parskip}\setlength{\parskip}{6pt}\setlength{\parindent}{0pt}
\usepackage{xcolor,tcolorbox,enumitem,listings,fancyhdr,titlesec,hyperref,booktabs,array}
\tcbuselibrary{skins,breakable}

\definecolor{deitelblue}{RGB}{0,70,127}
\definecolor{deitelgray}{RGB}{240,242,245}
\definecolor{deitelgreen}{RGB}{0,110,60}
\definecolor{deitellightblue}{RGB}{220,235,250}
\definecolor{deitelred}{RGB}{160,20,20}
\definecolor{deitellorange}{RGB}{200,90,0}

\newtcolorbox{boaprática}[1][]{colback=deitellightblue,colframe=deitelblue,
  fonttitle=\bfseries\small,title={Boa prática de programação},breakable,#1}
\newtcolorbox{erroco}[1][]{colback=red!5,colframe=deitelred,
  fonttitle=\bfseries\small,title={Erro comum de programação},breakable,#1}
\newtcolorbox{prev}[1][]{colback=yellow!8,colframe=deitellorange,
  fonttitle=\bfseries\small,title={Dica de prevenção de erro},breakable,#1}
\newtcolorbox{respostabox}[1][]{colback=deitelgray,colframe=deitelblue!60,
  fonttitle=\bfseries\small\color{deitelblue},title={Resposta detalhada},breakable,#1}
\newtcolorbox{enunciadoBox}[1][]{colback=white,colframe=deitelblue,
  fonttitle=\bfseries,title={#1},breakable}
\newtcolorbox{saida}[1][]{colback=black!88,colframe=deitelblue,
  fontupper=\ttfamily\small\color{green!70},breakable,#1}

\setlist[enumerate,1]{label=\textbf{\alph*)},leftmargin=2em}
\setlist[itemize]{leftmargin=2em}

\lstset{
  basicstyle=\ttfamily\small,backgroundcolor=\color{deitelgray},
  frame=single,framerule=0.4pt,rulecolor=\color{deitelblue!40},
  breaklines=true,showstringspaces=false,
  keywordstyle=\color{deitelblue}\bfseries,
  commentstyle=\color{deitelgreen}\itshape,
  stringstyle=\color{deitelred},
  language=C,numbers=left,numberstyle=\tiny\color{gray},
  xleftmargin=18pt,xrightmargin=5pt,stepnumber=1,
}

\pagestyle{fancy}\fancyhf{}
\fancyhead[L]{\small\color{deitelblue}\textbf{C: Como Programar -- 6ª ed. | Deitel \& Deitel}}
\fancyhead[R]{\small\color{deitelblue}Cap. 2 -- Autorrevisão}
\fancyfoot[C]{\small\thepage}
\renewcommand{\headrulewidth}{0.4pt}

\titleformat{\section}[block]{\large\bfseries\color{deitelblue}}{\thesection.}{0.5em}{}[\titlerule]
\titleformat{\subsection}[block]{\normalsize\bfseries\color{deitelblue!80}}{\thesubsection.}{0.5em}{}

\hypersetup{colorlinks=true,linkcolor=deitelblue,urlcolor=deitelblue}

\begin{document}

\begin{titlepage}
  \centering\vspace*{2cm}
  {\color{deitelblue}\rule{\linewidth}{2pt}}\\[0.4cm]
  {\huge\bfseries\color{deitelblue}C: Como Programar}\\[0.2cm]
  {\large\color{deitelblue}6ª Edição $\cdot$ Paul Deitel \& Harvey Deitel}\\[0.2cm]
  {\color{deitelblue}\rule{\linewidth}{2pt}}\\[1cm]
  {\LARGE\bfseries Capítulo 2}\\[0.3cm]
  {\Large\color{deitelblue!80}Introdução à Programação em C}\\[0.8cm]
  \begin{tcolorbox}[colback=deitellightblue,colframe=deitelblue,width=0.8\linewidth]
    \centering
    {\large\bfseries Exercícios de Autorrevisão}\\[0.2cm]
    {\normalsize Exercícios 2.1 e 2.2 -- Completamente resolvidos\\
    Abordagem incremental: Capítulos 1 e 2}
  \end{tcolorbox}
  \vfill{\small Abordagem incremental Deitel -- Caps. 1 e 2}
\end{titlepage}

\section*{Fundamentos Acumulados até o Capítulo~2}

Neste capítulo, você escreveu seus primeiros programas completos em C. Construindo sobre a
base do Capítulo~1 (hardware, software, linguagens e ambiente de desenvolvimento), o
Capítulo~2 introduz: a estrutura de um programa C (\texttt{main}, chaves, ponto e vírgula),
as funções \texttt{printf} e \texttt{scanf}, variáveis e tipos (\texttt{int}), operadores
aritméticos, precedência e a instrução \texttt{if}. Os exercícios de autorrevisão consolidam
esses fundamentos essenciais.

\begin{boaprática}
  Os conceitos do Capítulo~2 são a espinha dorsal de toda programação em C. Domine
  \texttt{printf}, \texttt{scanf}, variáveis, aritmética e \texttt{if} agora, pois eles
  serão usados em \textbf{todo} programa que você escreverá durante este curso.
\end{boaprática}

\tableofcontents\newpage

% ============================================================
\section{Exercício 2.1 -- Preenchimento de Lacunas}
% ============================================================

\begin{enunciadoBox}[Enunciado]
Preencha os espaços nas seguintes frases:
\begin{enumerate}
  \item Todo programa em C inicia sua execução pela função \underline{\hspace{3cm}}.
  \item O(a) \underline{\hspace{2cm}} inicia o corpo de cada função, e o(a) \underline{\hspace{2cm}} encerra o corpo de cada função.
  \item Toda instrução termina com um(a) \underline{\hspace{3cm}}.
  \item A função \underline{\hspace{2cm}} da biblioteca-padrão exibe informações na tela.
  \item A sequência de escape \texttt{\textbackslash n} representa o caractere de \underline{\hspace{2cm}}, que faz com que o cursor se posicione no início da próxima linha na tela.
  \item A função \underline{\hspace{2cm}} da biblioteca-padrão é usada para obter dados do teclado.
  \item O especificador de conversão \underline{\hspace{2cm}} é usado em uma string de controle de formato de \texttt{scanf} para indicar que um inteiro será digitado, e em uma string de controle de formato de \texttt{printf} para indicar que um inteiro será exibido.
  \item Sempre que um novo valor é colocado em uma posição da memória, ele anula o valor que ocupava essa mesma posição anteriormente. Esse processo é chamado de \underline{\hspace{3cm}}.
  \item Quando um valor é lido de uma posição da memória, o valor nessa posição é preservado; esse processo é chamado de \underline{\hspace{3cm}}.
  \item A instrução \underline{\hspace{2cm}} é usada na tomada de decisões.
\end{enumerate}
\end{enunciadoBox}

\bigskip

\subsection*{Item (a) -- A função que inicia a execução}

\begin{respostabox}
\textbf{Resposta:} \textcolor{deitelblue}{\textbf{\texttt{main}}}

\medskip\textbf{Explicação:}

Todo programa em C começa e organiza sua execução a partir da função \texttt{main}. Essa é
uma regra absoluta da linguagem: independentemente de quantas funções um programa tenha,
a execução sempre se inicia em \texttt{main}. A estrutura mínima de um programa C é:

\begin{lstlisting}
/* Estrutura minima de um programa C */
#include <stdio.h>

int main( void )
{
    /* instrucoes do programa */
    return 0;   /* indica conclusao com sucesso */
} /* fim da funcao main */
\end{lstlisting}

\begin{itemize}
  \item \texttt{int} $\to$ indica que \texttt{main} retorna um valor inteiro ao sistema operacional
  \item \texttt{void} $\to$ indica que \texttt{main} não recebe argumentos nessa forma simples
  \item \texttt{return 0;} $\to$ valor 0 sinaliza ao SO que o programa terminou com sucesso
\end{itemize}
\end{respostabox}

\subsection*{Item (b) -- Os delimitadores do corpo de uma função}

\begin{respostabox}
\textbf{Resposta:} \textcolor{deitelblue}{\textbf{chave à esquerda (\texttt{\{})} inicia / \textbf{chave à direita (\texttt{\}})} encerra}

\medskip\textbf{Explicação:}

Em C, o corpo de toda função é delimitado por um par de chaves. A \textbf{chave à esquerda}
\texttt{\{} marca o início do bloco de código da função, e a \textbf{chave à direita} \texttt{\}}
marca seu fim. Todo o código entre essas chaves compõe o \textit{corpo} da função.

\begin{lstlisting}
int main( void )
{           /* <- chave a esquerda: inicio do corpo */
    printf( "Bem-vindo a C!\n" );
    return 0;
}           /* <- chave a direita: fim do corpo */
\end{lstlisting}

O par de chaves com o código entre eles é chamado de \textbf{bloco}. Blocos serão
muito importantes ao estudarmos estruturas de controle nos Capítulos~3 e~4.
\end{respostabox}

\subsection*{Item (c) -- O terminador de instrução}

\begin{respostabox}
\textbf{Resposta:} \textcolor{deitelblue}{\textbf{ponto e vírgula (\texttt{;})}}

\medskip\textbf{Explicação:}

Em C, toda instrução executável deve ser encerrada com um \textbf{ponto e vírgula}, que é
chamado de \textit{terminador de instrução}. Exemplos:

\begin{lstlisting}
printf( "Ola, mundo!\n" );   /* instrucao de saida -- termina com ; */
int x;                        /* declaracao -- termina com ;          */
x = 5 + 3;                   /* instrucao de atribuicao -- termina com ; */
\end{lstlisting}

\begin{erroco}
  Esquecer o ponto e vírgula ao final de uma instrução é um dos erros de
  sintaxe mais comuns em C. O compilador emitirá uma mensagem de erro,
  geralmente na linha \textit{seguinte} à instrução onde o \texttt{;} está
  faltando, o que pode ser confuso.
\end{erroco}
\end{respostabox}

\subsection*{Item (d) -- A função que exibe informações na tela}

\begin{respostabox}
\textbf{Resposta:} \textcolor{deitelblue}{\textbf{\texttt{printf}}}

\medskip\textbf{Explicação:}

A função \texttt{printf} (\textit{print formatted}) da biblioteca-padrão \texttt{<stdio.h>}
é responsável por exibir informações formatadas na tela (no fluxo \texttt{stdout}).
Sua sintaxe geral é:

\begin{lstlisting}
printf( "string de controle de formato", arg1, arg2, ... );
\end{lstlisting}

Exemplos práticos:
\begin{lstlisting}
printf( "Bem-vindo a C!\n" );          /* texto simples         */
printf( "O valor e: %d\n", x );        /* exibe inteiro x       */
printf( "Soma = %d\n", a + b );        /* calcula e exibe       */
printf( "%d + %d = %d\n", a, b, a+b ); /* multiplos valores     */
\end{lstlisting}
\end{respostabox}

\subsection*{Item (e) -- A sequência de escape \texttt{\textbackslash n}}

\begin{respostabox}
\textbf{Resposta:} \textcolor{deitelblue}{\textbf{nova linha} (\textit{newline})}

\medskip\textbf{Explicação:}

A sequência de escape \texttt{\textbackslash n} representa o caractere de
\textbf{nova linha}. Quando \texttt{printf} encontra \texttt{\textbackslash n} em uma
string, move o cursor para o \textit{início da próxima linha} na tela.

\medskip
O caractere de barra invertida \texttt{\textbackslash} é chamado de
\textbf{caractere de escape}. Combinado com o próximo caractere, forma uma
\textbf{sequência de escape} com significado especial. Outras sequências úteis:

\begin{center}
\renewcommand{\arraystretch}{1.4}
\begin{tabular}{c l}
  \toprule \textbf{Sequência} & \textbf{Significado} \\ \midrule
  \texttt{\textbackslash n} & nova linha \\
  \texttt{\textbackslash t} & tabulação horizontal \\
  \texttt{\textbackslash\textbackslash} & barra invertida literal \\
  \texttt{\textbackslash "} & aspas duplas literais \\
  \bottomrule
\end{tabular}
\end{center}
\end{respostabox}

\subsection*{Item (f) -- A função que lê dados do teclado}

\begin{respostabox}
\textbf{Resposta:} \textcolor{deitelblue}{\textbf{\texttt{scanf}}}

\medskip\textbf{Explicação:}

A função \texttt{scanf} (\textit{scan formatted}) lê dados formatados da entrada-padrão
(\texttt{stdin}), que normalmente é o teclado. Sua sintaxe geral é:

\begin{lstlisting}
scanf( "string de controle de formato", &variavel1, &variavel2, ... );
\end{lstlisting}

O símbolo \texttt{\&} antes de cada variável é o \textbf{operador de endereço}, que
informa à \texttt{scanf} \textit{onde} na memória armazenar o valor lido. Exemplo:

\begin{lstlisting}
int idade;
printf( "Digite sua idade: " );
scanf( "%d", &idade );   /* le um inteiro e armazena em idade */
\end{lstlisting}

\begin{erroco}
  Esquecer o \texttt{\&} antes de uma variável em \texttt{scanf} é um dos erros
  mais frequentes de programadores iniciantes. Sem o \texttt{\&}, o programa
  pode causar uma ``falha de segmentação'' (acesso inválido à memória).
\end{erroco}
\end{respostabox}

\subsection*{Item (g) -- O especificador de conversão para inteiros}

\begin{respostabox}
\textbf{Resposta:} \textcolor{deitelblue}{\textbf{\texttt{\%d}}}

\medskip\textbf{Explicação:}

O especificador de conversão \texttt{\%d} é usado para indicar um valor inteiro decimal
(a letra \texttt{d} vem de \textit{decimal integer}). Ele funciona tanto em
\texttt{scanf} (para ler um inteiro) quanto em \texttt{printf} (para exibir um inteiro):

\begin{lstlisting}
int num;
scanf( "%d", &num );              /* le um inteiro do teclado */
printf( "Voce digitou: %d\n", num ); /* exibe o inteiro na tela  */
\end{lstlisting}

O \texttt{\%} em uma string de controle de formato sinaliza ao \texttt{scanf}/\texttt{printf}
que um \textbf{especificador de conversão} está começando. Tudo que vem após \texttt{\%} até
o caractere de formato especifica o tipo e o formato do dado.
\end{respostabox}

\subsection*{Item (h) -- Colocar um novo valor em uma posição de memória}

\begin{respostabox}
\textbf{Resposta:} \textcolor{deitelblue}{\textbf{destrutivo}}

\medskip\textbf{Explicação:}

Da Seção~2.4: quando um novo valor é colocado em uma posição de memória, ele
\textbf{substitui (destrói)} o valor que existia anteriormente naquela posição.
Por isso, essa operação é chamada de \textbf{destrutiva}.

\medskip Exemplo passo a passo:

\begin{lstlisting}
int x;
x = 10;   /* memoria: x = 10 */
x = 25;   /* memoria: x = 25 -- o valor 10 foi destruido! */
x = x + 1; /* memoria: x = 26 -- o valor 25 foi destruido! */
\end{lstlisting}

Na instrução \texttt{x = x + 1}, o computador primeiro \textit{lê} o valor de \texttt{x}
(operação não destrutiva), calcula \texttt{x + 1}, e então \textit{armazena} o resultado
em \texttt{x} (operação destrutiva -- o valor anterior é perdido).
\end{respostabox}

\subsection*{Item (i) -- Ler um valor sem alterar a memória}

\begin{respostabox}
\textbf{Resposta:} \textcolor{deitelblue}{\textbf{não destrutivo}}

\medskip\textbf{Explicação:}

Quando um valor é \textit{lido} de uma posição de memória para ser usado em um cálculo
ou exibição, o valor nessa posição é \textbf{preservado}. Por isso, a operação de leitura
é chamada de \textbf{não destrutiva}.

\medskip Exemplo:
\begin{lstlisting}
int a = 5, b = 3;
int soma = a + b;   /* le a (=5) e b (=3) -- nao destrutivo */
                    /* a ainda e 5, b ainda e 3 depois do calculo */
printf( "%d %d %d\n", a, b, soma ); /* imprime: 5 3 8 */
\end{lstlisting}

A distinção entre operações destrutivas e não destrutivas é fundamental para entender
o comportamento de variáveis e a instrução de atribuição em C.
\end{respostabox}

\subsection*{Item (j) -- A instrução de tomada de decisão}

\begin{respostabox}
\textbf{Resposta:} \textcolor{deitelblue}{\textbf{\texttt{if}}}

\medskip\textbf{Explicação:}

A instrução \texttt{if} (introduzida na Seção~2.6) permite que um programa tome decisões
com base na verdade ou falsidade de uma \textbf{condição}. Sua forma básica é:

\begin{lstlisting}
if ( condicao ) {
    /* executado apenas se condicao for verdadeira (nao zero) */
    printf( "A condicao e verdadeira!\n" );
} /* fim do if */
\end{lstlisting}

Condições são formadas com os \textbf{operadores relacionais} (\texttt{<}, \texttt{>},
\texttt{<=}, \texttt{>=}) e os \textbf{operadores de igualdade} (\texttt{==}, \texttt{!=}).

\begin{erroco}
  O erro mais clássico com \texttt{if} é confundir o operador de igualdade
  \texttt{==} com o operador de atribuição \texttt{=}. O código
  \texttt{if (x = 5)} \textit{não} verifica se \texttt{x} é igual a 5 --
  ele \textit{atribui} 5 a \texttt{x}, e a condição sempre será verdadeira!
  Leia \texttt{==} como ``dois iguais'' e \texttt{=} como ``recebe''.
\end{erroco}
\end{respostabox}

\newpage

% ============================================================
\section{Exercício 2.2 -- Verdadeiro ou Falso}
% ============================================================

\begin{enunciadoBox}[Enunciado]
Indique se cada uma das sentenças a seguir é \textbf{verdadeira} ou \textbf{falsa}.
Explique sua resposta no caso de alternativas \textit{falsas}.
\end{enunciadoBox}

\subsection*{Item (a) -- \texttt{printf} sempre começa no início de uma nova linha}

\begin{tcolorbox}[colback=red!5,colframe=deitelred,title={\bfseries FALSO}]
\textbf{Explicação:} A função \texttt{printf} começa a imprimir \textit{onde o cursor
estiver posicionado}, e isso pode ser em qualquer posição em uma linha. Por exemplo, se
você chamar \texttt{printf("Ola ")} e depois \texttt{printf("Mundo!\n")}, a saída será
\texttt{Ola Mundo!} tudo na mesma linha. Somente o caractere \texttt{\textbackslash n}
move o cursor para o início da próxima linha.
\end{tcolorbox}

\subsection*{Item (b) -- Comentários fazem o computador imprimir o texto na tela}

\begin{tcolorbox}[colback=red!5,colframe=deitelred,title={\bfseries FALSO}]
\textbf{Explicação:} Comentários em C (\texttt{/* ... */} ou \texttt{//}) são
completamente \textit{ignorados} pelo compilador -- eles não geram nenhum código de
máquina e não causam nenhuma ação durante a execução do programa. Sua única finalidade
é \textbf{documentar o código} e melhorar sua legibilidade para os humanos.

\begin{lstlisting}
/* Este texto NUNCA aparecera na tela -- e apenas um comentario */
printf( "Este texto SIM aparece na tela.\n" );
\end{lstlisting}
\end{tcolorbox}

\subsection*{Item (c) -- \texttt{\textbackslash n} move o cursor para o início da próxima linha}

\begin{tcolorbox}[colback=green!5,colframe=deitelgreen,title={\bfseries VERDADEIRO}]
\textbf{Explicação:} A sequência de escape \texttt{\textbackslash n} representa o caractere
de \textit{nova linha}. Quando o \texttt{printf} a encontra, posiciona o cursor no início
da próxima linha da tela. É exatamente isso que a torna tão essencial em todo programa C.
\end{tcolorbox}

\subsection*{Item (d) -- Todas as variáveis precisam ser declaradas antes de serem usadas}

\begin{tcolorbox}[colback=green!5,colframe=deitelgreen,title={\bfseries VERDADEIRO}]
\textbf{Explicação:} Em C padrão (ANSI C), toda variável deve ser declarada com seu nome
e tipo \textit{antes} de ser usada. Além disso, as declarações devem aparecer
\textbf{após a chave à esquerda de uma função e antes de quaisquer instruções executáveis}.
Tentar usar uma variável não declarada causa um erro de compilação.
\end{tcolorbox}

\subsection*{Item (e) -- Todas as variáveis precisam receber um tipo ao serem declaradas}

\begin{tcolorbox}[colback=green!5,colframe=deitelgreen,title={\bfseries VERDADEIRO}]
\textbf{Explicação:} Em C, toda declaração de variável especifica obrigatoriamente um
\textit{tipo de dado}. Neste capítulo, usamos o tipo \texttt{int} para inteiros.
Outros tipos que estudaremos mais adiante incluem \texttt{float}, \texttt{double} e
\texttt{char}. O tipo informa ao compilador quanto espaço de memória reservar e como
interpretar os bits armazenados naquela posição.
\end{tcolorbox}

\subsection*{Item (f) -- C considera \texttt{numero} e \texttt{NuMeRo} idênticas}

\begin{tcolorbox}[colback=red!5,colframe=deitelred,title={\bfseries FALSO}]
\textbf{Explicação:} C é uma linguagem \textbf{sensível a maiúsculas e minúsculas}
(\textit{case-sensitive}). Os identificadores \texttt{numero} e \texttt{NuMeRo} são
tratados como variáveis \textit{completamente diferentes}. Em C, cada letra conta:
\texttt{Main} não é o mesmo que \texttt{main}, por exemplo -- e isso causaria um
erro de compilação.
\begin{boaprática}
  Por convenção Deitel, nomes de variáveis simples devem começar com
  letra minúscula. Identificadores iniciados com maiúscula terão significado
  especial em C++ (classes), que estudaremos nos Capítulos 15--24.
\end{boaprática}
\end{tcolorbox}

\subsection*{Item (g) -- Declarações podem aparecer em qualquer parte do corpo de uma função}

\begin{tcolorbox}[colback=red!5,colframe=deitelred,title={\bfseries FALSO}]
\textbf{Explicação:} Em C padrão (ANSI C), as declarações de variáveis devem aparecer
\textbf{após a chave à esquerda que abre o corpo da função e antes de quaisquer instruções
executáveis}. Inserir uma declaração após uma instrução executável (como um \texttt{printf})
causa um \textit{erro de sintaxe} na maioria dos compiladores C padrão.

\begin{lstlisting}
int main( void )
{
    int a;               /* OK: declaracao antes das instrucoes */
    printf( "ok\n" );   /* instrucao executavel */
    int b;               /* ERRO em C padrao: declaracao apos instrucao */
    return 0;
}
\end{lstlisting}

\textit{Nota:} C99 e C++ permitem declarações em qualquer ponto, mas neste livro
seguimos o padrão C clássico.
\end{tcolorbox}

\subsection*{Item (h) -- Argumentos de \texttt{printf} precisam ser precedidos por \texttt{\&}}

\begin{tcolorbox}[colback=red!5,colframe=deitelred,title={\bfseries FALSO}]
\textbf{Explicação:} O operador de endereço \texttt{\&} é obrigatório nos argumentos de
\texttt{scanf} (para que ela saiba \textit{onde} armazenar o valor lido), mas
\textbf{não} deve ser usado nos argumentos de \texttt{printf} (que apenas \textit{lê}
o valor da variável para exibi-lo).

\begin{lstlisting}
int x = 42;
printf( "%d\n", x );    /* CORRETO: sem & em printf */
printf( "%d\n", &x );   /* ERRO: & em printf e incorreto */
scanf( "%d", &x );      /* CORRETO: com & em scanf       */
scanf( "%d", x );       /* ERRO: sem & em scanf e perigoso */
\end{lstlisting}

As exceções a essa regra (relacionadas a arrays e ponteiros) serão discutidas nos
Capítulos~6 e~7.
\end{tcolorbox}

\subsection*{Item (i) -- O operador módulo só pode ser usado com operandos inteiros}

\begin{tcolorbox}[colback=green!5,colframe=deitelgreen,title={\bfseries VERDADEIRO}]
\textbf{Explicação:} O operador módulo \texttt{\%} produz o resto da divisão inteira e
só é válido com operandos do tipo \texttt{int} (e outros tipos inteiros). Tentar usá-lo
com números de ponto flutuante (\texttt{float}, \texttt{double}) causa um erro de
compilação. Exemplo:
\begin{lstlisting}
int a = 17, b = 5;
printf( "Resto: %d\n", a % b ); /* Resultado: 2 -- CORRETO */
/* double x = 17.0 % 5.0; -- ERRO: % nao funciona com double */
\end{lstlisting}
\end{tcolorbox}

\subsection*{Item (j) -- Os operadores \texttt{*}, \texttt{/}, \texttt{\%}, \texttt{+} e \texttt{-} têm o mesmo nível de precedência}

\begin{tcolorbox}[colback=red!5,colframe=deitelred,title={\bfseries FALSO}]
\textbf{Explicação:} Em C, os operadores aritméticos possuem \textbf{dois} níveis de
precedência distintos (Seção~2.5):

\begin{center}
\renewcommand{\arraystretch}{1.4}
\begin{tabular}{c l l}
  \toprule \textbf{Nível} & \textbf{Operadores} & \textbf{Avaliação} \\ \midrule
  Maior & \texttt{*} \quad \texttt{/} \quad \texttt{\%} & da esquerda para a direita \\
  Menor & \texttt{+} \quad \texttt{-} & da esquerda para a direita \\ \bottomrule
\end{tabular}
\end{center}

Portanto, \texttt{2 + 3 * 4} é avaliado como \texttt{2 + 12 = 14} (não como \texttt{5 * 4 = 20}),
porque \texttt{*} tem maior precedência que \texttt{+}.
\end{tcolorbox}

\subsection*{Item (k) -- Nomes com mais de 31 caracteres são idênticos em todos os sistemas C}

\begin{tcolorbox}[colback=red!5,colframe=deitelred,title={\bfseries FALSO}]
\textbf{Explicação:} O padrão C garante que \textbf{apenas os primeiros 31 caracteres}
de um identificador precisam ser reconhecidos pelos compiladores. Portanto, dois
identificadores que diferem apenas após o 31º caractere podem ser tratados como
\textit{idênticos} por alguns compiladores, enquanto outros com suporte maior podem
distingui-los. Os nomes do enunciado \texttt{esteéumnúmerosuperhiperlongo1234567} e
\texttt{esteéumnúmerosuperhiperlongo1234568} diferem apenas no 44º caractere -- isso
pode causar problemas de portabilidade.

\begin{boaprática}
  Use identificadores com no máximo 31 caracteres para garantir portabilidade
  entre todos os compiladores C padrão.
\end{boaprática}
\end{tcolorbox}

\subsection*{Item (l) -- Um programa que exibe três linhas precisa de três \texttt{printf}}

\begin{tcolorbox}[colback=red!5,colframe=deitelred,title={\bfseries FALSO}]
\textbf{Explicação:} Uma única instrução \texttt{printf} pode exibir múltiplas linhas
usando várias sequências de escape \texttt{\textbackslash n}:

\begin{lstlisting}
/* Uma instrucao printf exibindo tres linhas */
printf( "Linha 1\nLinha 2\nLinha 3\n" );
\end{lstlisting}

\begin{saida}
Linha 1\\
Linha 2\\
Linha 3
\end{saida}

Cada \texttt{\textbackslash n} na string de controle de formato gera um salto de linha.
\end{tcolorbox}

\newpage
\section*{Gabarito Rápido -- Capítulo 2: Autorrevisão}

\begin{tcolorbox}[colback=deitellightblue,colframe=deitelblue,title={\bfseries\large Respostas Resumidas}]
\textbf{Exercício 2.1:}
\begin{enumerate}[label=\alph*)]
  \item \texttt{main}
  \item chave à esquerda (\texttt{\{}); chave à direita (\texttt{\}})
  \item ponto e vírgula (\texttt{;})
  \item \texttt{printf}
  \item nova linha (\textit{newline})
  \item \texttt{scanf}
  \item \texttt{\%d}
  \item destrutivo
  \item não destrutivo
  \item \texttt{if}
\end{enumerate}

\medskip\textbf{Exercício 2.2 -- Verdadeiro/Falso:}
\begin{itemize}[noitemsep]
  \item[a)] \textcolor{deitelred}{\textbf{Falso}} -- \texttt{printf} imprime onde o cursor estiver
  \item[b)] \textcolor{deitelred}{\textbf{Falso}} -- comentários são ignorados pelo compilador
  \item[c)] \textcolor{deitelgreen}{\textbf{Verdadeiro}}
  \item[d)] \textcolor{deitelgreen}{\textbf{Verdadeiro}}
  \item[e)] \textcolor{deitelgreen}{\textbf{Verdadeiro}}
  \item[f)] \textcolor{deitelred}{\textbf{Falso}} -- C é sensível a maiúsculas e minúsculas
  \item[g)] \textcolor{deitelred}{\textbf{Falso}} -- declarações devem vir antes das instruções executáveis
  \item[h)] \textcolor{deitelred}{\textbf{Falso}} -- \texttt{\&} é usado em \texttt{scanf}, não em \texttt{printf}
  \item[i)] \textcolor{deitelgreen}{\textbf{Verdadeiro}}
  \item[j)] \textcolor{deitelred}{\textbf{Falso}} -- \texttt{*,/,\%} têm precedência maior que \texttt{+,-}
  \item[k)] \textcolor{deitelred}{\textbf{Falso}} -- pode haver diferença após os 31 primeiros caracteres
  \item[l)] \textcolor{deitelred}{\textbf{Falso}} -- um único \texttt{printf} pode exibir várias linhas
\end{itemize}
\end{tcolorbox}

\end{document}
